% GENERATED FILE. Edit canonical lessons, then run npm run generate:books.
% canonical-source-hash: fnv1a64:f7cd70327ce8f91f
% canonical-lessons: KA-C395-adalita, KA-C395-sacivalaya, KA-C395-pancayiti, KA-C395-jarigolisu, KA-C395-samiti

\chapter{Running Public Offices}
\label{ch:ka-running-public-offices}
\hlchaptermodality{395}

\begin{chapteropening}
\textbf{By the end of this chapter:} \emph{I can say administration, a ministry, a village council, to implement and a committee.}

Complete the last lesson of chapter 395: I can say administration, a ministry, a village council, to implement and a committee.
\end{chapteropening}

\section[\texorpdfstring{\kn{ಆಡಳಿತ} (āḍaḷita)}{ಆಡಳಿತ (āḍaḷita)}]{\kn{ಆಡಳಿತ} (āḍaḷita) — administration}

\label{lesson:KA-C395-adalita}



\begin{quote}
Before the new one: say the Kannada for to boycott, then the Kannada for leadership.
\end{quote}

\subsection*{You'll want to know: \kn{ಆಡಳಿತ}}
\textbf{\kn{ಆಡಳಿತ}} — \emph{āḍaḷita} — \textquotedblleft{}administration\textquotedblright{}.

\begin{quote}
\kn{ಊರಿನ} \kn{ಆಡಳಿತ} \kn{ಚೆನ್ನಾಗಿ} \kn{ನಡೆಯುತ್ತಿದೆ}.
\end{quote}

\begin{quote}
\emph{ūrina āḍaḷita cennāgi naḍeyuttide.}
\end{quote}

\begin{quote}
The town's administration is running well.
\end{quote}

\kn{ಆಡಳಿತ} \kn{ಪಕ್ಷ}, the ruling party.

\begin{cousinweb}
\textbf{\kn{ಆಡಳಿತ}} (\emph{āḍaḷita}) is usually traced to the native Dravidian \emph{āḷ}, \textquotedblleft{}to possess, to rule, to manage\textquotedblright{} (DEDR 5157); DEDR does not list \emph{āḍaḷita} itself, so how the word was formed is not settled.
\end{cousinweb}

\begin{grammarlens}[title={What you've built}]
One more word for giving reasons and opinions.
\end{grammarlens}

\subsection*{Your turn}
\begin{itemize}
\raggedright
  \item \emph{Say it:} \emph{āḍaḷita}
  \item \emph{Say it:} \emph{āḍaḷita}, once more
  \item \emph{Say it:} \emph{ūrina āḍaḷita cennāgi naḍeyuttide.} — The town's administration is running well.
  \item \emph{Recall:} say the Kannada for to boycott, then the Kannada for leadership, then say \emph{āḍaḷita} again
\end{itemize}

\begin{tcolorbox}[breakable,colback=teal!4,colframe=teal!35!black,title={Before you move on}]
What does \kn{ಆಡಳಿತ} mean? (\textquotedblleft{}Administration\textquotedblright{}.) Say it once more.
\end{tcolorbox}
\section[\texorpdfstring{\kn{ಸಚಿವಾಲಯ} (sacivālaya)}{ಸಚಿವಾಲಯ (sacivālaya)}]{\kn{ಸಚಿವಾಲಯ} (sacivālaya) — a ministry}

\label{lesson:KA-C395-sacivalaya}



\begin{quote}
Before the new one: say the Kannada for leadership, then the Kannada for administration.
\end{quote}

\subsection*{You'll want to know: \kn{ಸಚಿವಾಲಯ}}
\textbf{\kn{ಸಚಿವಾಲಯ}} — \emph{sacivālaya} — \textquotedblleft{}a ministry\textquotedblright{}.

\begin{quote}
\kn{ಅವನ} \kn{ಅಪ್ಪ} \kn{ಸಚಿವಾಲಯದಲ್ಲಿ} \kn{ಕೆಲಸ} \kn{ಮಾಡುತ್ತಾರೆ}.
\end{quote}

\begin{quote}
\emph{avana appa sacivālayadalli kelasa māḍuttāre.}
\end{quote}

\begin{quote}
His father works in the ministry.
\end{quote}

\kn{ಸಚಿವ} (minister) + \kn{ಆಲಯ} (house).

\begin{cousinweb}
\textbf{\kn{ಸಚಿವಾಲಯ}} (\emph{sacivālaya}) is built from Sanskrit \emph{saciva}, \textquotedblleft{}a companion, a king's minister\textquotedblright{} (Monier-Williams), and \emph{ālaya}, \textquotedblleft{}a house\textquotedblright{} (CDIAL 1366).
\end{cousinweb}

\begin{grammarlens}[title={What you've built}]
One more word for giving reasons and opinions.
\end{grammarlens}

\subsection*{Your turn}
\begin{itemize}
\raggedright
  \item \emph{Say it:} \emph{sacivālaya}
  \item \emph{Say it:} \emph{sacivālaya}, once more
  \item \emph{Say it:} \emph{avana appa sacivālayadalli kelasa māḍuttāre.} — His father works in the ministry.
  \item \emph{Recall:} say the Kannada for leadership, then the Kannada for administration, then say \emph{sacivālaya} again
\end{itemize}

\begin{tcolorbox}[breakable,colback=teal!4,colframe=teal!35!black,title={Before you move on}]
What does \kn{ಸಚಿವಾಲಯ} mean? (\textquotedblleft{}A ministry\textquotedblright{}.) Say it once more.
\end{tcolorbox}
\section[\texorpdfstring{\kn{ಪಂಚಾಯಿತಿ} (pañcāyiti)}{ಪಂಚಾಯಿತಿ (pañcāyiti)}]{\kn{ಪಂಚಾಯಿತಿ} (pañcāyiti) — a village council}

\label{lesson:KA-C395-pancayiti}



\begin{quote}
Before the new one: say the Kannada for administration, then the Kannada for a ministry.
\end{quote}

\subsection*{You'll want to know: \kn{ಪಂಚಾಯಿತಿ}}
\textbf{\kn{ಪಂಚಾಯಿತಿ}} — \emph{pañcāyiti} — \textquotedblleft{}a village council\textquotedblright{}.

\begin{quote}
\kn{ಪಂಚಾಯಿತಿ} \kn{ಊರಿನಲ್ಲಿ} \kn{ಹೊಸ} \kn{ಶಾಲೆ} \kn{ಕಟ್ಟಿಸಿತು}.
\end{quote}

\begin{quote}
\emph{pañcāyiti ūrinalli hosa śāle kaṭṭisitu.}
\end{quote}

\begin{quote}
The village council had a new school built.
\end{quote}

\kn{ಪಂಚಾಯಿತಿ} \kn{ಚುನಾವಣೆ}, a council election.

\begin{cousinweb}
\textbf{\kn{ಪಂಚಾಯಿತಿ}} (\emph{pañcāyiti}) is a loan from Hindustani \emph{pañcāyat}, \textquotedblleft{}a council\textquotedblright{}, once a body of five elders, from Sanskrit \emph{pañca}, \textquotedblleft{}five\textquotedblright{} (CDIAL 7655).
\end{cousinweb}

\begin{grammarlens}[title={What you've built}]
One more word for giving reasons and opinions.
\end{grammarlens}

\subsection*{Your turn}
\begin{itemize}
\raggedright
  \item \emph{Say it:} \emph{pañcāyiti}
  \item \emph{Say it:} \emph{pañcāyiti}, once more
  \item \emph{Say it:} \emph{pañcāyiti ūrinalli hosa śāle kaṭṭisitu.} — The village council had a new school built.
  \item \emph{Recall:} say the Kannada for administration, then the Kannada for a ministry, then say \emph{pañcāyiti} again
\end{itemize}

\begin{tcolorbox}[breakable,colback=teal!4,colframe=teal!35!black,title={Before you move on}]
What does \kn{ಪಂಚಾಯಿತಿ} mean? (\textquotedblleft{}A village council\textquotedblright{}.) Say it once more.
\end{tcolorbox}
\section[\texorpdfstring{\kn{ಜಾರಿಗೊಳಿಸು} (jārigoḷisu)}{ಜಾರಿಗೊಳಿಸು (jārigoḷisu)}]{\kn{ಜಾರಿಗೊಳಿಸು} (jārigoḷisu) — to implement}

\label{lesson:KA-C395-jarigolisu}



\begin{quote}
Before the new one: say the Kannada for a ministry, then the Kannada for a village council.
\end{quote}

\subsection*{You'll want to know: \kn{ಜಾರಿಗೊಳಿಸು}}
\textbf{\kn{ಜಾರಿಗೊಳಿಸು}} — \emph{jārigoḷisu} — \textquotedblleft{}to implement\textquotedblright{}.

\begin{quote}
\kn{ಸರ್ಕಾರ} \kn{ಹೊಸ} \kn{ಯೋಜನೆಯನ್ನು} \kn{ಜಾರಿಗೊಳಿಸಿತು}.
\end{quote}

\begin{quote}
\emph{sarkāra hosa yōjaneyannu jārigoḷisitu.}
\end{quote}

\begin{quote}
The government implemented the new scheme.
\end{quote}

Past (I …): \kn{ಜಾರಿಗೊಳಿಸಿದೆ}, \emph{jārigoḷiside}.

\kn{ಜಾರಿಗೆ} \kn{ಬರು}, to come into force.

\begin{cousinweb}
\textbf{\kn{ಜಾರಿಗೊಳಿಸು}} (\emph{jārigoḷisu}) is built on \emph{jārī}, \textquotedblleft{}current, in force\textquotedblright{}, a loan from Hindustani that comes from Arabic \emph{jārī}, \textquotedblleft{}flowing\textquotedblright{}, with the Kannada ending -\emph{goḷisu}, \textquotedblleft{}to make\textquotedblright{}.
\end{cousinweb}

\begin{grammarlens}[title={What you've built}]
One more word for giving reasons and opinions.
\end{grammarlens}

\subsection*{Your turn}
\begin{itemize}
\raggedright
  \item \emph{Say it:} \emph{jārigoḷisu}
  \item \emph{Say it:} \emph{jārigoḷisu}, once more
  \item \emph{Say it:} \emph{sarkāra hosa yōjaneyannu jārigoḷisitu.} — The government implemented the new scheme.
  \item \emph{Recall:} say the Kannada for a ministry, then the Kannada for a village council, then say \emph{jārigoḷisu} again
\end{itemize}

\begin{tcolorbox}[breakable,colback=teal!4,colframe=teal!35!black,title={Before you move on}]
What does \kn{ಜಾರಿಗೊಳಿಸು} mean? (\textquotedblleft{}To implement\textquotedblright{}.) Say it once more.
\end{tcolorbox}
\section[\texorpdfstring{\kn{ಸಮಿತಿ} (samiti)}{ಸಮಿತಿ (samiti)}]{\kn{ಸಮಿತಿ} (samiti) — a committee}

\label{lesson:KA-C395-samiti}



\begin{quote}
Before the new one: say the Kannada for a village council, then the Kannada for to implement.
\end{quote}

\subsection*{You'll want to know: \kn{ಸಮಿತಿ}}
\textbf{\kn{ಸಮಿತಿ}} — \emph{samiti} — \textquotedblleft{}a committee\textquotedblright{}.

\begin{quote}
\kn{ಹಬ್ಬದ} \kn{ಕೆಲಸಕ್ಕೆ} \kn{ಒಂದು} \kn{ಸಮಿತಿ} \kn{ಮಾಡಿದರು}.
\end{quote}

\begin{quote}
\emph{habbada kelasakke ondu samiti māḍidaru.}
\end{quote}

\begin{quote}
They formed a committee for the festival work.
\end{quote}

\kn{ಸಮಿತಿಯ} \kn{ಅಧ್ಯಕ್ಷ}, the chair of a committee.

\begin{cousinweb}
\textbf{\kn{ಸಮಿತಿ}} (\emph{samiti}) is borrowed from Sanskrit \emph{samiti}, \textquotedblleft{}a meeting, an association\textquotedblright{} (Monier-Williams).
\end{cousinweb}

\begin{grammarlens}[title={What you've built}]
One more word for giving reasons and opinions.
\end{grammarlens}

\subsection*{Your turn}
\begin{itemize}
\raggedright
  \item \emph{Say it:} \emph{samiti}
  \item \emph{Say it:} \emph{samiti}, once more
  \item \emph{Say it:} \emph{habbada kelasakke ondu samiti māḍidaru.} — They formed a committee for the festival work.
  \item \emph{Recall:} say the Kannada for a village council, then the Kannada for to implement, then say \emph{samiti} again
\end{itemize}

\begin{tcolorbox}[breakable,colback=teal!4,colframe=teal!35!black,title={Before you move on}]
What does \kn{ಸಮಿತಿ} mean? (\textquotedblleft{}A committee\textquotedblright{}.) Say it once more.
\end{tcolorbox}
